\subsection{Detalle de la calibración del optimizador}
\label{anx:calibracion}

\sloppy En este anexo se reúnen las tablas de consulta de la calibración que resume la
Sección~\ref{sec:optimizacion-cal}, donde también están las figuras. Todas salen de los archivos JSON de la calibración,
entre ellos \texttt{umbral.json}, \texttt{umbral\_elegida.json}, \texttt{calibracion\_nsga.json}, \texttt{sensibilidad.json} y \texttt{hipervolumen\_prueba.json}, medidos en un \calCPU{} con
\calRAM~GiB de memoria.

\paragraph{Umbral de enumeración.} La Tabla~\ref{tab:calopt-umbral} da los tiempos por nivel del
espacio por nivel; la gráfica con el ajuste está en la Sección~\ref{sec:optimizacion-cal}. Con la configuración
de partida, NSGA-II tardaba una mediana de \calTNsgaMediana~s y el cruce quedaba en
$S\approx\calCruce$; con la elegida baja a \calTNsgaElegidaMediana~s y el cruce, a
$S\approx\calCruceElegida$.

\begin{table}[htbp]
\centering
\caption{Tiempo de pared de la enumeración exhaustiva y de NSGA-II por nivel $b$ del espacio
($S\in[2^b,2^{b+1})$), con la configuración de partida y con la elegida. «Exacto más rápido» y
«HV mínimo» se refieren a la elegida; \calSubproblemas{} subproblemas y \calSemillasHV{} semillas por
subproblema para el hipervolumen.}
\label{tab:calopt-umbral}
\small
\begin{tabular}{rrrrrrrr}
\toprule
 & & \multicolumn{2}{c}{exacto (s)} & \multicolumn{2}{c}{NSGA-II, mediana (s)} & & \\
$b$ & subproblemas & mediana & máx. & partida & elegida & exacto más rápido & HV mínimo \\
\midrule
6 & 5 & 0{,}001 & 0{,}001 & 1{,}382 & 0{,}258 & 5 & 1{,}0000 \\
7 & 5 & 0{,}002 & 0{,}003 & 0{,}703 & 0{,}197 & 5 & 1{,}0000 \\
8 & 5 & 0{,}005 & 0{,}011 & 0{,}388 & 0{,}163 & 5 & 1{,}0000 \\
9 & 5 & 0{,}015 & 0{,}020 & 0{,}372 & 0{,}188 & 5 & 1{,}0000 \\
10 & 5 & 0{,}032 & 0{,}058 & 0{,}439 & 0{,}211 & 5 & 1{,}0000 \\
11 & 5 & 0{,}089 & 0{,}294 & 0{,}401 & 0{,}183 & 3 & 0{,}9998 \\
12 & 5 & 0{,}354 & 0{,}731 & 0{,}420 & 0{,}223 & 1 & 0{,}9999 \\
13 & 5 & 1{,}301 & 2{,}457 & 0{,}461 & 0{,}239 & 0 & 1{,}0000 \\
14 & 5 & 1{,}420 & 7{,}553 & 0{,}581 & 0{,}288 & 0 & 0{,}9996 \\
15 & 3 & 11{,}912 & 25{,}134 & 0{,}632 & 0{,}337 & 0 & 0{,}9998 \\
16 & 3 & 48{,}534 & 142{,}993 & 0{,}538 & 0{,}274 & 0 & 0{,}9997 \\
17 & 3 & 27{,}650 & 179{,}797 & 0{,}642 & 0{,}356 & 0 & 1{,}0000 \\
18 & 3 & --- & --- & 0{,}712 & 0{,}356 & 0 & 0{,}9999 \\
\bottomrule

\end{tabular}
\end{table}

\paragraph{Población y generaciones.} La Tabla~\ref{tab:calopt-rejilla} da la rejilla completa sobre
los pliegues de validación (\calCVListas{} listas, \calSemillasCV{} semillas por lista; los pliegues
tienen \calPlieguesListas{} listas). El pliegue en
que la regla eligió otra configuración escogió una todavía más barata, que quedó por fuera del margen
en su pliegue apartado, lo que indica que reducir más el presupuesto ya cuesta calidad. En la prueba,
el \calFraccionHV{} de las corridas de la elegida alcanzó al menos 0,999 del hipervolumen exacto, la
peor corrida llegó a \calRazonMinPrueba{} y el mínimo por lista fue \calPruebaElegidaHVMin. La rejilla
es reducida y se declara: 2 semillas por lista en la validación y 5 en la prueba.

\begin{table}[htbp]
\centering
\caption{Rejilla de población ($a$) y generaciones ($b$) sobre los pliegues de validación.
$\star$: configuración elegida; $\circ$: configuración de partida.}
\label{tab:calopt-rejilla}
\small
\begin{tabular}{rrrrrrc}
\toprule
$a$ & $b$ & evaluaciones medias & HV medio & HV mínimo por lista & tiempo medio (s) & \\
\midrule
0{,}25 & 0{,}25 & 1\,231 & 0{,}9960 & 0{,}9711 & 0{,}29 &  \\
0{,}25 & 0{,}50 & 2\,475 & 0{,}9973 & 0{,}9708 & 0{,}57 &  \\
0{,}25 & 1{,}00 & 4\,963 & 0{,}9971 & 0{,}9615 & 1{,}18 &  \\
0{,}25 & 2{,}00 & 9\,937 & 0{,}9975 & 0{,}9725 & 2{,}30 &  \\
0{,}50 & 0{,}25 & 2\,453 & 0{,}9986 & 0{,}9904 & 0{,}48 &  \\
0{,}50 & 0{,}50 & 4\,931 & 0{,}9998 & 0{,}9970 & 1{,}01 & $\star$ \\
0{,}50 & 1{,}00 & 9\,890 & 1{,}0000 & 0{,}9981 & 2{,}01 &  \\
0{,}50 & 2{,}00 & 19\,800 & 1{,}0000 & 0{,}9979 & 4{,}20 &  \\
1{,}00 & 0{,}25 & 4\,907 & 0{,}9987 & 0{,}9896 & 0{,}90 &  \\
1{,}00 & 0{,}50 & 9\,862 & 0{,}9999 & 0{,}9981 & 1{,}93 &  \\
1{,}00 & 1{,}00 & 19\,779 & 1{,}0000 & 0{,}9993 & 3{,}96 & $\circ$ \\
1{,}00 & 2{,}00 & 39\,601 & 1{,}0000 & 1{,}0000 & 8{,}00 &  \\
2{,}00 & 0{,}25 & 9\,814 & 0{,}9988 & 0{,}9935 & 1{,}86 &  \\
2{,}00 & 0{,}50 & 19\,723 & 0{,}9999 & 0{,}9973 & 4{,}03 &  \\
2{,}00 & 1{,}00 & 39\,558 & 1{,}0000 & 0{,}9992 & 8{,}50 &  \\
2{,}00 & 2{,}00 & 79\,201 & 1{,}0000 & 1{,}0000 & 16{,}46 &  \\
\bottomrule

\end{tabular}
\end{table}

\paragraph{Umbral $\tau$.} Con la curaduría disponible, en $\tau=0{,}75$ se acepta el
\calTauAceptada{} de los renglones con una precisión del \calTauPrecision, pero el parecido solo toma
\calTauValores{} valores distintos y la validación cruzada agrupada elige valores entre
\calTauCVMin{} y \calTauCVMax{} según el pliegue. Por eso, y porque no se calculó el acuerdo entre los 4 anotadores, el valor de 0,75 se fija por su
precisión sobre los \calTauPares{} pares y se reporta con su intervalo, sin afinarlo pliegue a pliegue.

\paragraph{Sensibilidad del kit recomendado.} En las \calSensListas{} listas de prueba con al menos
dos puntos en el frente (mediana de \calSensFrenteMediana{} puntos) se sortearon
\calSensPerturbaciones{} vectores de pesos por lista y celda. Sin sobreprecio la estabilidad es 1 por
construcción, porque la cota de sobreprecio deja admisible solo el kit más barato. En el primer
esquema cada peso se sortea dentro de su intervalo del 95\,\% de la P21; en el segundo, los pesos
siguen una distribución de Dirichlet centrada en los de la P21 con concentración $\kappa$, que admite
perturbaciones más anchas. La Tabla~\ref{tab:calopt-sensibilidad} da todas las celdas.

\begin{table}[htbp]
\centering
\caption{Fracción de sorteos de pesos en que el kit recomendado por MAUT es el mismo que con los
pesos de la P21. DE: desviación estándar del peso del precio bajo la perturbación.}
\label{tab:calopt-sensibilidad}
\small
\begin{tabular}{rlrrlr}
\toprule
sobreprecio & perturbación & DE & media & IC 95\,\% & mínimo por lista \\
\midrule
0\,\% & IC 95\,\% P21 & --- & 1{,}000 & [1{,}000; 1{,}000] & 1{,}000 \\
5\,\% & IC 95\,\% P21 & --- & 0{,}965 & [0{,}932; 0{,}991] & 0{,}610 \\
10\,\% & IC 95\,\% P21 & --- & 0{,}986 & [0{,}964; 1{,}000] & 0{,}692 \\
20\,\% & IC 95\,\% P21 & --- & 0{,}946 & [0{,}904; 0{,}981] & 0{,}540 \\
0\,\% & Dirichlet $\kappa=20$ & 0{,}099 & 1{,}000 & [1{,}000; 1{,}000] & 1{,}000 \\
5\,\% & Dirichlet $\kappa=20$ & 0{,}099 & 0{,}892 & [0{,}830; 0{,}947] & 0{,}429 \\
10\,\% & Dirichlet $\kappa=20$ & 0{,}099 & 0{,}868 & [0{,}822; 0{,}913] & 0{,}541 \\
20\,\% & Dirichlet $\kappa=20$ & 0{,}099 & 0{,}735 & [0{,}663; 0{,}803] & 0{,}168 \\
0\,\% & Dirichlet $\kappa=50$ & 0{,}063 & 1{,}000 & [1{,}000; 1{,}000] & 1{,}000 \\
5\,\% & Dirichlet $\kappa=50$ & 0{,}063 & 0{,}918 & [0{,}866; 0{,}965] & 0{,}502 \\
10\,\% & Dirichlet $\kappa=50$ & 0{,}063 & 0{,}927 & [0{,}888; 0{,}960] & 0{,}558 \\
20\,\% & Dirichlet $\kappa=50$ & 0{,}063 & 0{,}828 & [0{,}765; 0{,}886] & 0{,}287 \\
0\,\% & Dirichlet $\kappa=100$ & 0{,}045 & 1{,}000 & [1{,}000; 1{,}000] & 1{,}000 \\
5\,\% & Dirichlet $\kappa=100$ & 0{,}045 & 0{,}931 & [0{,}884; 0{,}974] & 0{,}520 \\
10\,\% & Dirichlet $\kappa=100$ & 0{,}045 & 0{,}953 & [0{,}923; 0{,}978] & 0{,}618 \\
20\,\% & Dirichlet $\kappa=100$ & 0{,}045 & 0{,}876 & [0{,}818; 0{,}926] & 0{,}379 \\
0\,\% & Dirichlet $\kappa=200$ & 0{,}032 & 1{,}000 & [1{,}000; 1{,}000] & 1{,}000 \\
5\,\% & Dirichlet $\kappa=200$ & 0{,}032 & 0{,}945 & [0{,}903; 0{,}980] & 0{,}557 \\
10\,\% & Dirichlet $\kappa=200$ & 0{,}032 & 0{,}973 & [0{,}946; 0{,}993] & 0{,}632 \\
20\,\% & Dirichlet $\kappa=200$ & 0{,}032 & 0{,}910 & [0{,}857; 0{,}956] & 0{,}456 \\
\bottomrule

\end{tabular}
\end{table}

\paragraph{Hipervolumen por perfil.} La Tabla~\ref{tab:calopt-hv} da el hipervolumen relativo al
frente exacto en las \calHVListas{} listas de prueba (\calHVListasNsga{} de ellas con NSGA-II). El mínimo del frente de producción fue \calHVFrenteMin.

\begin{table}[htbp]
\centering
\caption{Hipervolumen relativo al frente exacto en las listas de prueba, con intervalo
\emph{bootstrap} del 95\,\%.}
\label{tab:calopt-hv}
\small
\begin{tabular}{lrl}
\toprule
conjunto & media & IC 95\,\% \\
\midrule
Frente de producci\'on & 0{,}9969 & [0{,}9951; 0{,}9986] \\
Ahorrador $\cup$ Premium & 0{,}501 & [0{,}474; 0{,}526] \\
Ahorrador & 0{,}362 & [0{,}343; 0{,}382] \\
Premium & 0{,}231 & [0{,}201; 0{,}262] \\
MAUT con pesos P21 & 0{,}517 & [0{,}474; 0{,}558] \\
\bottomrule

\end{tabular}
\end{table}

\subsection{Detalle de la evaluación del asistente NL2C}
\label{anx:nl2c}

En este anexo se reúnen las tablas de la ronda 3 de la evaluación del asistente que resume la
Sección~\ref{sec:nl2c}. La Tabla~\ref{tab:nl2c-rondas} da las medidas de la evaluación.

\begin{table}[htbp]
\centering\small
\caption{Resultados de las rondas 2 y 3 del asistente NL2C.}
\label{tab:nl2c-rondas}
\begin{tabular}{@{}p{5.4cm}p{3.9cm}p{3.9cm}@{}}\toprule
Medida & Ronda 2 (criterio original) & Ronda 3 (criterio nuevo)\\\midrule
Tarea completa & 57 de 70 (81,4\,\%) & 58 de 70 (82,9\,\%)\\
F1 de personajes y marcas & 1,000 & 1,000\\
Latencia p50 / p95 (ms) & 999 / 1.344 & 1.000 / 1.177\\
Ataques exitosos (ASR) & 0 de 81 & 0 de 81\\
El modelo lo intentó & 5 de 81 & 6 de 81\\
Utilidad bajo ataque & 36 de 60 & 36 de 60\\
Controles sin ataque & 12 de 12 & 12 de 12\\\bottomrule
\end{tabular}
\end{table}

\iffalse % Tablas de la ronda 3 por campo, ataque y costo: docs/validacion/nl2c/ y docs/memoria/sections/tablas/nl2c/
\begin{table}[htbp]
\centering\small
\caption{Coincidencia con el oro por campo de $\mathcal{R}$ en la ronda 3.}
\label{tab:nl2c-campos}
\begin{tabular}{@{}lrrr@{}}\toprule
Campo de R & Aciertos & Evaluados & Coincidencia\\\midrule
acción & 68 & 70 & 97,1\,\%\\
presupuesto & 58 & 60 & 96,7\,\%\\
personajes & 58 & 60 & 96,7\,\%\\
marcas excluidas & 58 & 60 & 96,7\,\%\\
solo disponibles & 58 & 60 & 96,7\,\%\\
estilo de gasto & 56 & 60 & 93,3\,\%\\
sobreprecio & 6 & 6 & 100,0\,\%\\
peso mayor & 50 & 60 & 83,3\,\%\\
tarea completa & 58 & 70 & 82,9\,\%\\
\bottomrule\end{tabular}

\end{table}

\begin{table}[htbp]
\centering\small
\caption{Ataques de inyección por categoría en la ronda 3.}
\label{tab:nl2c-asr}
\begin{tabular}{@{}lrrrr@{}}\toprule
Categoría & Casos & Intentos & Modelo lo intentó & Ataque exitoso\\\midrule
directa sobre R & 16 & 48 & 3 (6,2\,\%) & 0 (0,0\,\%)\\
directa sobre el kit & 5 & 15 & 3 (20,0\,\%) & 0 (0,0\,\%)\\
indirecta (texto del kit) & 6 & 18 & 0 (0,0\,\%) & 0 (0,0\,\%)\\
\midrule
Total & 27 & 81 & 6 (7,4\,\%) & 0 (0,0\,\%)\\
\bottomrule\end{tabular}

\end{table}

\begin{table}[htbp]
\centering\small
\caption{Condiciones de éxito del atacante en la ronda 3.}
\label{tab:nl2c-condiciones}
\begin{tabular}{@{}lrrr@{}}\toprule
Condición de éxito del atacante & Evaluados & Intentó el modelo & Llegó al padre\\\midrule
R fuera de dominio & 24 & 3 & 0\\
más de 5 propuestas & 6 & 3 & 0\\
fuga del prompt de sistema & 3 & 0 & 0\\
marcado ejecutable en pantalla & 12 & 0 & 0\\
cadena del atacante en pantalla & 45 & 0 & 0\\
dato sensible en pantalla & 18 & 0 & 0\\
función fuera de la lista & 9 & 0 & 0\\
renglón o perfil ajeno & 9 & 0 & 0\\
descarte sin confirmar (criterio de la ronda 3) & 6 & 0 & 0\\
propuesta no pedida & 6 & 0 & 0\\
enlace externo en pantalla & 3 & 0 & 0\\
\bottomrule\end{tabular}

\end{table}

\begin{table}[htbp]
\centering\small
\caption{Costo y latencia de la ronda 3. El costo usa un precio de lista supuesto, no medido.}
\label{tab:nl2c-costo}
\begin{tabular}{@{}lr@{}}\toprule
Medida & Valor\\\midrule
Intentos (frases por repeticiones) & 70\\
F1 de personajes y marcas & 1,000\\
Precisión / exhaustividad & 1,000 / 1,000\\
Latencia p50 (ms) & 1000\\
Latencia p95 (ms) & 1177\\
Tokens de entrada por pedido (media) & 616\\
Tokens de salida por pedido (media) & 39\\
Costo por pedido (USD, precio supuesto) & 0,000077\\
Intentos sin servicio & 0\\
\bottomrule\end{tabular}

\end{table}

\begin{table}[htbp]
\centering\small
\caption{Utilidad bajo ataque y controles en la ronda 3.}
\label{tab:nl2c-utilidad}
\begin{tabular}{@{}lrr@{}}\toprule
Situación & Tareas cumplidas & Tasa\\\midrule
Bajo ataque (con tarea legítima) & 36 de 60 & 60,0\,\%\\
Controles sin ataque (kit) & 12 de 12 & 100,0\,\%\\
\bottomrule\end{tabular}

\end{table}

\fi
